\chapter{Gráficos}\label{chp:graficos}

    O \cref{grf:distidadeingr} a seguir traz a distribuição dos alunos ingressantes no Curso de Ciências da Computação por faixa etária.

    \begin{graphic}[htb]
        \caption{Barras}
        \label{grf:distidadeingr}
        \centering
        \begin{tikzpicture}
            \tiny
            \begin{axis}[
                %title = {Distribuição dos Alunos por Faixa Etária},
                %xlabel = {Períodos letivos},
                ylabel = {Qtd. de alunos},
                font =\footnotesize,
                width = 10cm,
                height = 4.9cm,
                ybar =2*\pgflinewidth,
                bar width=8pt,
                ymajorgrids = true,
                symbolic x coords={$2013.1$,$2013.2$,$2014.1$,$2014.2$},
                xtick = data,
                ymin=0,
                ymax=50,
                enlarge x limits=0.25,
                legend cell align=left,
                legend image code/.code={\draw[#1] (0cm,-0.1cm) rectangle (0.2cm,0.1cm); },
                legend style={legend pos=outer north east, column sep=1ex}
                %legend style={at={(1.0,1.15)}, anchor=north,legend columns=-1}
            ]
    
                \addplot[style={cl00,fill=cl01}] coordinates {($2013.1$, 27) ($2013.2$, 30) ($2014.1$, 35) ($2014.2$, 44)};
                \addplot[style={cl00,fill=cl02}] coordinates {($2013.1$, 11) ($2013.2$,  8) ($2014.1$,  5) ($2014.2$,  6)};
                \addplot[style={cl00,fill=cl03}] coordinates {($2013.1$,  0) ($2013.2$,  0) ($2014.1$,  0) ($2014.2$,  1)};
                \addplot[style={cl00,fill=cl04}] coordinates {($2013.1$,  0) ($2013.2$,  1) ($2014.1$,  0) ($2014.2$,  0)};
    
                \legend{$16$ a $20$ anos, $21$ a $25$ anos, $26$ a $30$ anos, Acima de $30$ anos}
            \end{axis}
        \end{tikzpicture}
        \fonte{Autoria própria}
    \end{graphic}

    O \cref{grf:distidadematr} a seguir traz a distribuição dos alunos do Curso de Ciências da Computação, matriculados nos dois últimos anos, por faixa etária.

    \begin{graphic}[htb]
        \caption{Pizza}
        \label{grf:distidadematr}
        \centering
        \begin{tikzpicture}
            \tiny
            \pie[text=legend, color = {cl01,cl02,cl03,cl04}, radius=2]
                {38.62/16 a 20, 39.77/21 a 25, 14.99/26 a 30, 6.62/mais de 30}
        \end{tikzpicture}
        \fonte{Autoria própria}
    \end{graphic}

    O \cref{grf:funcoes} a seguir traz exemplos de plotagem de funções.

    \begin{graphic}[!htb]
        \caption{Funções}
        \label{grf:funcoes}
        \centering
        \begin{tikzpicture}
            \tiny
            \begin{axis}[
                axis y line = center,
                axis x line = middle,
                width = 10cm,
    %           height = 6cm,
                xlabel = $x$,
                ylabel = {$f(x)$},
                legend style={legend pos=outer north east, column sep=1ex}
            ]
    
                \addplot [domain =-20:20, samples=100, color=red ]{x^2 - 9};
                \addlegendentry {$f(x) = x^2 - 9$}
    
                \addplot [domain =-20:20, samples=100, color=blue]{x + 2};
                \addlegendentry {$g(x) = x + 2$}
            \end{axis}
        \end{tikzpicture}
        \fonte{Autoria própria}
    \end{graphic}

    \FloatBarrier

    O \cref{grf:fxetariaanual} traz uma representação em linhas da distribuição dos alunos matriculados, por ano e faixa etária, no Curso de Ciências da Computação.

    \begin{graphic}[!htb]
        \caption{Linhas}
        \label{grf:fxetariaanual}
        \centering
        \tiny
        \begin{tikzpicture}
            \begin{axis}[
                title = {Distribuição dos Alunos por Faixa Etária Anual},
                xlabel = {Anos},
                ylabel = {Qtd. de Alunos},
                font = \footnotesize,
                width = 0.6\textwidth,
                height = 6cm,
                ymajorgrids = true,
                grid style={dotted},
                x tick label style={/pgf/number format/1000 sep=},
                xtick = data,
                ytick = {0,40,80,...,200},
                ymin = 0,
                ymax = 200,
                enlarge x limits = 0.25,
                legend cell align = left,
                legend style={legend pos=outer north east, column sep=1ex}
            ]
                \addplot[mark=square*,line width=2pt,style={cl01}] coordinates {(2011,56) (2012,77) (2013,102) (2014,134)};
                \addlegendentry{16 a 20}
    
                \addplot[mark=o,line width=2pt,style={cl02}] coordinates {(2011,170) (2012,151) (2013,142) (2014,138)};
                \addlegendentry{21 a 25}
    
                \addplot[mark=triangle*,line width=2pt,style={cl03}] coordinates {(2011,73) (2012,66) (2013,55) (2014,52)};
                \addlegendentry{26 a 30}
    
                \addplot[mark=halfsquare*,line width=2pt,style={cl04}] coordinates {(2011,37) (2012,30) (2013,28) (2014,23)};
                \addlegendentry{Acima de 30}
            \end{axis}
        \end{tikzpicture}
        \fonte{Autoria própria}
    \end{graphic}